这是一道通信题。在这里描述发送端（Alice）与接收端（Bob）各自已知的信息、需要完成的任务以及相关约束。

你只需要提交一份完整程序，评测系统会对每组测试数据将其独立运行两次：

\begin{itemize}
\item 第一次运行（\texttt{first}）：扮演发送端，读入第一轮输入并输出通信消息。
\item 第二次运行（\texttt{second}）：扮演接收端，读入评测系统转发的通信消息与附加输入，输出最终答案。
\end{itemize}

两次运行为完全独立的进程，全局变量、内存与本地文件均不会保留，只能通过第一次运行的标准输出传递信息。

每次运行的输入第一行均为字符串 \texttt{first} 或 \texttt{second}（不含引号）。请先读入该标记进入对应分支，完成本轮输出后立即结束程序。

\InputFile

\textbf{第一次运行（\texttt{first}）}

第一行包含一个字符串 \texttt{first}。

在这里说明第一轮输入格式、参数含义及数据范围。

\textbf{第二次运行（\texttt{second}）}

第一行包含一个字符串 \texttt{second}。

在这里说明第二轮收到的通信消息格式、附加输入及数据范围。

\OutputFile

\textbf{第一次运行（\texttt{first}）}

在这里说明通信消息的输出格式、长度上限、字符集或数值范围限制。

\textbf{第二次运行（\texttt{second}）}

在这里说明最终答案的输出格式。

\Note

两次运行各自独立享有完整的时间与内存限制。

请不要输出任何额外提示语或调试信息，也不要在第一次运行结束后继续等待第二次输入。

在这里补充样例说明或本地测试方法（如有）。
